(1) 每种载流子横向电导均有$1/(1+\mu_j^2B^2)$因子，而电子与空穴的Hall项符号相反。取$\mathbf J=\begin{pmatrix}\sigma_{xx}&\sigma_{xy}\\-\sigma_{xy}&\sigma_{xx}\end{pmatrix}\mathbf E$，得$\sigma_{xx}=e[n\mu_e/(1+\mu_e^2B^2)+p\mu_h/(1+\mu_h^2B^2)]$、$\sigma_{xy}=eB[p\mu_h^2/(1+\mu_h^2B^2)-n\mu_e^2/(1+\mu_e^2B^2)]$。
(2) 对二维反对称矩阵直接求逆，$\rho=\sigma^{-1}=(\sigma_{xx}^2+\sigma_{xy}^2)^{-1}\begin{pmatrix}\sigma_{xx}&-\sigma_{xy}\\\sigma_{xy}&\sigma_{xx}\end{pmatrix}$。在$J_y=0$的Hall测量条件下，$E_y/J_x=\rho_{yx}=\sigma_{xy}/(\sigma_{xx}^2+\sigma_{xy}^2)$，纵向电阻率$\rho_{xx}=\sigma_{xx}/(\sigma_{xx}^2+\sigma_{xy}^2)$，不能单独用$1/\sigma_{xx}$替代。
(3) 在$|\mu_jB|\ll1$下，$\sigma_{xx}=e(n\mu_e+p\mu_h)+O(B^2)$，$\sigma_{xy}=eB(p\mu_h^2-n\mu_e^2)+O(B^3)$。因此$R_H=\lim_{B\to0}\rho_{yx}/B=(p\mu_h^2-n\mu_e^2)/[e(n\mu_e+p\mu_h)^2]$。Hall电压反号不必等于$n=p$，而是取决于带迁移率平方权重的补偿条件$p\mu_h^2=n\mu_e^2$。
(4) 若$n=p=N$且两迁移率均为$\mu$，则$\sigma_{xy}=0$严格对任意$B$成立，而$\sigma_{xx}=2eN\mu/(1+\mu^2B^2)$。因此$\rho_{xx}(B)=(1+\mu^2B^2)/(2eN\mu)$，相对磁电阻$[\rho_{xx}(B)-\rho_{xx}(0)]/\rho_{xx}(0)=\mu^2B^2$，在理想模型中无饱和。真实材料的多带散射、密度不均匀及量子振荡可改变这一高场结果。
